% 09 — VSS：从零到快照读路径与多卷闭包

\section{起点：活扫描读不了被锁定的文件}

文件级备份默认 \texttt{CreateFile + ReadFile}。
SQL Server、Outlook、甚至正在写入的日志常持有 **exclusive share** → \winapi{ERROR\_SHARING\_VIOLATION}，scan issue \texttt{locked}。

\begin{evolbox}[GAP-CONSIST-01 / Phase 14–18]
\textbf{Phase 14 (ABI v31)}：\texttt{VssSession} — COM + \winapi{IVssBackupComponents} 基础生命周期。\\
\textbf{Phase 15 (ABI v32)}：crash/app/auto 模式、多卷 junction 闭包、路径映射。\\
\textbf{Phase 18 (ABI v35)}：影子存储 WMI 预检、\texttt{--vss-fallback-live}。
\end{evolbox}

\textbf{定位}：VSS 为\textbf{文件级备份}提供一致性读路径；\textbf{不做}块级 raw 镜像（与 Veeam 整机备份不同）。

\section{一致性模式：crash / app / auto}

\begin{longtable}{@{}llp{6.5cm}@{}}
\caption{VSS 模式对照} \\
\toprule
模式 & C API flag & 行为 \\
\midrule
crash & \texttt{EB\_BACKUP\_FLAG\_VSS} & \winapi{DoSnapshotSet} 后直接读 shadow；无 writer 协调 \\
app & \texttt{+ EB\_BACKUP\_FLAG\_VSS\_APP} & \winapi{GatherWriterMetadata} → 快照 → 读 → \winapi{GatherWriterStatus} → \winapi{BackupComplete} \\
auto & \texttt{set\_vss\_mode("auto")} & 尝试 app；writer 失败降级 crash，issue \texttt{vss\_writer\_degraded} \\
\bottomrule
\end{longtable}

\begin{designbox}[为何读完后才 BackupComplete]
Writer 协议要求：快照后、释放 writer 前完成数据复制。
ebbackup 在 \textbf{scan + chunk 全部完成} 后调用 \texttt{FinishBackup}，避免 writer 超时与 torn read。
\end{designbox}

\section{VSS 生命周期}

\begin{figure}[htbp]
\centering
\begin{tikzpicture}[node distance=0.65cm]
  \node[wbox] (init) {InitializeForBackup};
  \node[wbox, below=of init] (meta) {GatherWriterMetadata（app 模式）};
  \node[wbox, below=of meta] (b) {StartSnapshotSet + AddToSnapshotSet};
  \node[wbox, below=of b] (snap) {DoSnapshotSet};
  \node[wbox, below=of snap] (read) {ScanDirectory + Enrich + Chunk（shadow 路径）};
  \node[wbox, below=of read] (fin) {GatherWriterStatus + BackupComplete};
  \node[wbox, below=of fin] (end) {EndBackup / Delete Snapshots};
  \draw[warr] (init) -- (meta) -- (b) -- (snap) -- (read) -- (fin) -- (end);
\end{tikzpicture}
\caption{VSS 生命周期与读顺序}
\end{figure}

实现：\filepath{vss\_session.cc}，封装 \texttt{VssSession::Begin / FinishBackup / End}。

\section{Begin：闭包、预检、快照}

\begin{lstlisting}[caption={VssSession::Begin 核心 (vss\_session.cc)}]
CheckPrerequisites();  // Admin 或 SeBackupPrivilege

ComputeVolumeClosure(logical_roots, closure_opts, &closure);
info_.cross_volume = closure.size() > 1;

if (!opts.skip_shadow_preflight) {
  CheckShadowStoragePreflightEx(closure, opts.shadow_storage_min_bytes, &storage);
  // 不足 → 失败或 fallback live（--vss-fallback-live）
}

impl_->StartSnapshotSet();
for (const auto& spec : closure)
  impl_->AddToSnapshotSet(spec, &entry.snapshot_id);
impl_->DoSnapshotSet();

BuildPathMap(mount_prefix → shadow_prefix);
\end{lstlisting}

\section{权限与特权}

推荐**管理员**或持有：
\begin{itemize}[nosep]
  \item \winapi{SeBackupPrivilege} — 读 bypass ACL；
  \item \winapi{SeRestorePrivilege} — 配合 restore 场景。
\end{itemize}

非 elevated：\filepath{vss\_locked\_file\_backup\_test} 可能 \texttt{GTEST\_SKIP}。

\section{多卷闭包：vss\_volume\_closure}

用户备份 \texttt{C:/project}，其中 junction 指向 \texttt{D:/data}。
\textbf{扫描不展开} D:（专册 06），但 VSS 需对 D: **做快照**才能一致读 junction target 上的文件（若用户另行包含该路径）。

\begin{lstlisting}[caption={ComputeVolumeClosure（概念）}]
ResolveVolumeForPath(root) → volume_key (\\?\Volume{GUID}\);
ProbeJunctionVolumes(root, depth=junction_probe_depth, &extra_volumes);
closure = unique(volume_keys);
\end{lstlisting}

\begin{lstlisting}[caption={ProbeJunctionVolumes BFS}]
// 深度 junction_probe_depth（默认 2）
// IO_REPARSE_TAG_MOUNT_POINT → ReadJunctionTargetUtf8
// 跨卷 target → 加入 closure，不 recurse 扫描
\end{lstlisting}

报告：\texttt{vss\_cross\_volume}、\texttt{vss\_volumes[]}（卷 GUID 列表）。

\section{影子存储预检：vss\_shadow\_storage}

\begin{enumerate}
  \item \textbf{WMI 优先}：\winapi{Win32\_ShadowStorage} 查询 Used/Max Space；
  \item \textbf{Fallback}：解析 \texttt{vssadmin list shadowstorage} 文本输出；
  \item \texttt{CheckShadowStoragePreflightEx}：默认 **512MB** 余量（\texttt{shadow\_storage\_min\_bytes}）；
  \item 失败 + \texttt{--vss-fallback-live}：降级活扫描，issue \texttt{vss\_shadow\_storage\_low}。
\end{enumerate}

CLI：\texttt{eb vss status}；JSON 报告 \texttt{vss\_shadow\_storage\_bytes[]}、\texttt{vss\_shadow\_storage\_ok}。

\section{路径映射：vss\_path\_map}

快照后逻辑源路径前缀映射到 shadow 设备：

\begin{phasebox}[路径映射示例]
源：\texttt{C:/data/file.txt}\\
Shadow：\texttt{\textbackslash?\textbackslash?\textbackslash Volume\{xxx\}\textbackslash data\textbackslash file.txt}\\
\texttt{ScanDirectory} 的 \texttt{walk\_root} 使用 shadow 前缀；manifest \texttt{relative\_path} 仍为 \texttt{data/file.txt}（对用户透明）。
\end{phasebox}

\texttt{MapToShadow} / \texttt{MapFromShadow} 在 session 内维护双向表。

\section{与 Win 元数据 enrichment 的衔接}

VSS 不改变 enrichment 代码路径——仅在 \textbf{映射后}绝对路径上调用 \texttt{CaptureWinMetaFromPath}：

\begin{center}
\begin{tabular}{@{}ll@{}}
\toprule
板块 & VSS 交互 \\
\midrule
02 扫描 enrichment & shadow 路径上 ACL/inode/reparse \\
04 ADS & \texttt{FindFirstStreamW} 对快照一致 \\
05 硬链 & 快照时间点 inode 稳定 \\
06 Junction & 闭包扩卷；scan 仍不展开 \\
07 Sparse & IOCTL run map 在 shadow 有效 \\
08 EFS & 加密属性在快照上可读 \\
\bottomrule
\end{tabular}
\end{center}

\section{备份报告字段（ABI v32+）}

JSON report 扩展：
\begin{itemize}[nosep]
  \item \texttt{vss\_used}、\texttt{vss\_mode}、\texttt{vss\_consistency}；
  \item \texttt{vss\_writers[]}、\texttt{vss\_snapshot\_set\_id}；
  \item \texttt{vss\_shadow\_storage\_ok}、\texttt{vss\_shadow\_storage\_bytes[]}；
  \item issue：\texttt{vss\_writer\_degraded}、\texttt{vss\_shadow\_storage\_low}。
\end{itemize}

\section{Job 扩展（Phase 18）}

\texttt{quiesce\_profile}、\texttt{vss\_app\_failure\_policy} 映射 Writer 失败策略。
\textbf{不含} SQL/Exchange 专用 Writer 插件——依赖系统已注册 VSS writer。

\section{明确不做}

\begin{itemize}[nosep]
  \item 块级镜像 / 裸金属恢复；
  \item 影子空间管理 GUI；
  \item Hyper-V/VMware 集成（Veeam B\&R 级 VM 备份）。
\end{itemize}

\section{测试矩阵}

\begin{longtable}{@{}lp{8.5cm}@{}}
\toprule
测试文件 & 覆盖 \\
\midrule
\filepath{vss\_path\_map\_test.cc} & 前缀映射 roundtrip \\
\filepath{vss\_lifecycle\_test.cc} & Begin/Finish/End COM 生命周期 \\
\filepath{vss\_volume\_closure\_test.cc} & junction 卷发现 \\
\filepath{vss\_shadow\_storage\_test.cc} & WMI/vssadmin 解析 \\
\filepath{vss\_locked\_file\_backup\_test.cc} & elevated 锁定文件可读 \\
\filepath{vss\_session\_stub\_test.cc} & 无 VSS 环境 stub \\
\bottomrule
\end{longtable}

\section{从零搭建 VSS 读路径的设计教训}

\begin{enumerate}
  \item \textbf{先 crash 后 app}：crash 模式已解决 80\% 锁定文件；app 模式增加 writer 失败面，需 auto 降级。
  \item \textbf{闭包与 scan 解耦}：用户 scope 由 scan 决定；VSS 卷集由闭包决定——两者故意不一致。
  \item \textbf{预检优于运行时失败}：影子存储不足时在 Begin 阶段 fail 或 fallback，避免扫到一半快照被删。
  \item \textbf{读完成再 BackupComplete}：与 Microsoft VSS 文档顺序一致，减少 SQL/AD writer 报错。
\end{enumerate}

\section{源码索引}

\begin{longtable}{@{}lp{9cm}@{}}
\toprule
文件 & 职责 \\
\midrule
\filepath{src/winmeta/vss\_session.cc} & COM 会话、快照、路径映射、FinishBackup \\
\filepath{src/winmeta/vss\_volume\_closure.cc} & 卷解析、junction BFS 闭包 \\
\filepath{src/winmeta/vss\_shadow\_storage.cc} & WMI/vssadmin 预检 \\
\filepath{src/scan/scan\_entry.cc} & walk\_root 使用 shadow 映射 \\
\filepath{include/ebbackup/winmeta/vss\_session.h} & \texttt{VssBeginOptions}、报告 DTO \\
\bottomrule
\end{longtable}
